\chapter{Análisis de viabilidad}
\label{chap:viabilidade}

\lettrine{E}{n este capítulo} se analiza la viabilidad del proyecto desde tres perspectivas complementarias: dimensión temporal, dimensión tecnológica y dimensión económica. 

\section{Viabilidad temporal}
Para la realización de este TFM se ha estimado una duración aproximada de cuatro meses. Dicha planificación incluye las fases de análisis, diseño, implementación, pruebas y redacción de la memoria. Además, se contempla la familiarización de las nuevas tecnologías. Dado el alcance definido la viabilidad temporal es aceptable y no representa un riesgo crítico para la finalización del proyecto.

La metodología usada utiliza desarrollos cortos para ir validando el progreso y ajustando el alcance si fuera necesario.

\section{Viabilidad tecnológica}
Para el estudio de viablidad tecnológica se han considerado las tecnologías heredadas del TFG, así como las nuevas tecnologías propuestas para el desarrollo de la aplicación móvil y la integración de notificaciones push.

En resumen, el stack tecnológico propuesto para el TFM es el siguiente:
\begin{itemize}
  \item Backend: Spring Boot (Java) para servicios REST y la integración con PostgreSQL/PostGIS.
  \item Base de datos: PostgreSQL con PostGIS para soporte espacial.
  \item Frontend web: React con Material-UI para interfaces ricas.
  \item Aplicación Android: Kotlin y Jetpack Compose para la capa de presentación; cliente nativo para acceder a sensores y recibir notificaciones push.
  \item Notificaciones push: Firebase Cloud Messaging (FCM) para la entrega fiable de mensajes a dispositivos Android.
\end{itemize}

Kotlin y FCM merecen una mención específica: ambas tecnologías están ampliamente adoptadas, cuentan con extensa documentación y comunidades activas, y existen guías oficiales para su integración con backends Java/Spring. Por tanto, su adopción no supone un riesgo tecnológico significativo; aportan ventajas en productividad, interoperabilidad móvil y soporte nativo para notificaciones.

Las demás tecnologías son conocidas por el autor gracias a su dedicación profesional y al proyecto previo, lo que reduce el riesgo de adopción y acelera el desarrollo.


\section{Viabilidad económica}
El proyecto no requiere una inversión económica significativa en hardware o licencias de software. El software utilizado es mayoritariamente Open Source (Spring Boot, PostgreSQL, React, Kotlin/Jetpack Compose), por lo que no hay costes de licencia en el entorno de desarrollo académico.

En el ámbito del desarrollo y las pruebas estos pueden ejecutarse con recursos disponibles por lo que no es necesario adquirir nuevo hardware.Sobre el uso de servicios en la nube o APIs de terceros(FMC, MapLibre, OpenWeatherMap) pueden surgir costes, pero se pueden controlar gracias a los planes gratuitos o a la monitorización de uso.

La dimensión económica no supone un riesgo significativo para la ejecución del TFM pues los costes directos son bajos.

\section{Conclusión del análisis de viabilidad}
En resumen, el proyecto es totalmente viable porque combina un plan de trabajo realista de cuatro meses con un uso inteligente de la tecnología. Al reutilizar la base del TFG y apoyarse en herramientas gratuitas y potentes como Kotlin y FCM, el riesgo técnico es mínimo y el coste es prácticamente cero. Esta experiencia previa, sumada a una planificación clara, garantiza que el proyecto se pueda terminar a tiempo y con la calidad necesaria para un TFM.

